\uuid{GPOD}
\titre{Famille libre/Génératrice/Base/Polynômes}
\auteur{hanani ; mzari}
\organisation{exo7}
\licence{CC BY-NC-SA 4.0}
\datecreate{2020-01-20}
\niveau{L1}
\module{Algèbre}
\chapitre{Espace vectoriel}
\sousChapitre{Base}
\difficulte{2}
\theme{Base et dimension, Famille libre, Génératrice, Base, Polynômes}
\exo7qcmid{lille-2.espaces-vectoriels.032}

\contenu{
\begin{question}
\qclassification{106.04}
\qtags{motcle=Famille libre/Génératrice/Base/Polynômes}
Dans $\Rr_2[X]$, l'ensemble des polynômes à coefficients réels de degré $\le 2$, on considère les polynômes $P_1= X, P_2= X(X+1), P_3 = (X+1)^2$. Quelles sont les assertions vraies ?
\begin{answers}
\good{$\{P_1,P_2,P_3 \}$ est une famille libre}.
\bad{$\{P_1+P_2,P_3 \}$ est une famille génératrice de $\Rr_2[X]$}.
\good{$\{P_1,P_2,P_3 \}$ est une base de $\Rr_2[X]$}.
\bad{$\{P_2,P_3 \}$ est une base de $\Rr_2[X]$}.
\end{answers}
\begin{explanations} On vérifie que $\{P_1,P_2,P_3 \}$ est une famille libre de $\Rr_2[X]$. De plus, cette famille contient
$3$ polynômes  et la dimension de $\Rr_2[X]$ est $3$, donc c'est une base de $\Rr_2[X]$.
\end{explanations}
\end{question}
}
